\chapter*{本册结构}
\addcontentsline{toc}{chapter}{本册结构}

第三册不按“AI 名词清单”组织，而按一台 C++20 本地大模型推理引擎的建设顺序组织。它明确承接第一册和第二册：第一册给出 C++、二进制、汇编、ABI、对象布局和可信测量；第二册给出现代 CPU、cache/TLB、SIMD、多线程、NUMA、运行时、I/O 和背压模型；第三册把这些能力集中用在模型、张量、AI 编译器、推理 runtime 和 CPU/GPU 后端上。

读者即使完全没有学过编译原理，也应该先拿到全书地图：输入是什么，语义在哪里被检查，中间表示怎样建立，优化改变了哪条机器账本，后端怎样落到 CPU/GPU，运行时怎样证明结果没有跑偏。

\begin{center}
\begin{minipage}{0.92\linewidth}
\begin{verbatim}
model files / tokenizer / weights
  -> manifest and schema resolution
  -> shape, dtype and quantization verifier
  -> quantization calibration, int8/int4 and KV cache policy
  -> graph IR / tensor IR / loop IR
  -> fusion, memory planning and lowering
  -> CPU backend / GPU backend executable plans
  -> runtime session, KV cache, scheduler and sampler
  -> oracle, profiler, benchmark and framework comparison
\end{verbatim}
\end{minipage}
\end{center}

\section*{纵向样板：MiniLLM-CPU 从 toy 到 production slice}

第三册不再把“toy demo”和“生产系统”割裂。贯穿样板叫 MiniLLM-CPU，它分三层尺度推进：

\begin{center}
\begin{tabular}{p{0.18\linewidth}p{0.36\linewidth}p{0.30\linewidth}}
\toprule
尺度 & 固定内容 & 训练目标 \\
\midrule
Tiny 数值层 & \code{vocab=16, H=8, L=2, heads=2, kv\_heads=1, seq=4} & 手算 shape、stride、KV offset、logits golden \\
Kernel 工程层 & \code{T=2,K=8,N=6} 和 \code{H=4096,K=4096,O=4096} & scalar、packing、SIMD、tail、perf/objdump \\
7B 系统层 & \code{L=32,H=4096,heads=32,kv\_heads=8,ctx=4096} & 权重字节、KV 容量、tokens/s 上限、SLO \\
\bottomrule
\end{tabular}
\end{center}

这三层不是三个项目。Tiny 层定义正确性，kernel 层训练热路径，7B 层训练工程判断。每个章节都要说明自己在哪一层留下证据。例如 RoPE/KV cache 先用 tiny 层手算 position 和 offset，再用 7B 层计算每 token KV 字节；GEMV 先用 \code{T=2,K=8,N=6} 画 packed layout，再用 \code{4096} 级 shape 做带宽上限和 benchmark；serving 先用 tiny session 验证取消和事件顺序，再用 7B KV 预算推 admission。

\section*{固定端到端 trace}

第三册最重要的链路不是目录，而是一条 token trace：

\begin{lstlisting}[numbers=none,caption={MiniLLM-CPU 固定端到端 trace}]
prompt
  -> tokenizer ids
  -> embedding [T,H]
  -> RMSNorm [T,H]
  -> Q/K/V projection
  -> RoPE(Q,K,position)
  -> KV cache append/read
  -> causal attention with GQA
  -> residual + RMSNorm
  -> SwiGLU MLP
  -> lm_head logits
  -> sampler
  -> next token
\end{lstlisting}

每一步都要能回答七个问题：

\begin{enumerate}
  \item shape 是什么。
  \item dtype 是什么。
  \item stride 和物理 layout 是什么。
  \item C++ 对象是谁拥有这块内存。
  \item reference 输出或 oracle checkpoint 在哪里。
  \item 优化后端怎样改变地址流或执行计划。
  \item trace/report 怎样证明这一步没有漂移、没有偷换后端、没有把成本算错。
\end{enumerate}

如果某章不能把自己放进这条 trace，就应删减或重写。AI Infra 不是名词覆盖面，而是同一个 token 从业务输入一路走到机器执行、再走回可解释输出。

\section*{AI compiler 的硬交付}

第四章不能只说 manifest、IR、pass、lowering。它必须产出可检查文件：

\begin{lstlisting}[numbers=none,caption={AI compiler 必须导出的 artifact}]
manifest-report.json:
  tensors, tokenizer contract, dtype, quant profile

ir-before.json:
  imported graph, tensor values, state edges

ir-after-passes.json:
  canonical graph, fusion, view, liveness facts

memory-plan.json:
  arena size, offsets, workspace, KV cache binding

executable-plan.json:
  prefill/decode segments, backend kernels, profiler labels
\end{lstlisting}

这些 artifact 要覆盖正例和负例。正例证明 tiny 模型能编译成 CPU plan；负例证明坏 embedding 行数、坏 QKV shape、坏 \code{kv\_heads}、坏 quant group、fusion 删除 oracle checkpoint 都能被准确拒绝。若 compiler 只在内存里悄悄做事，读者无法学习 pass pipeline，也无法在面试或作品集中展示工程证据。

\section*{CPU microkernel 的硬交付}

第五章 CPU 部分的最低样板是 int8 linear/GEMV。它不能只写“用 SIMD 优化”，必须形成完整证据链：

\begin{lstlisting}[numbers=none,caption={CPU int8 linear 样板交付}]
1. fp32 reference GEMV
2. int8 scalar baseline
3. packed int8 layout and roundtrip oracle
4. AVX2/AVX-512 kernel path
5. tail path for K/N not divisible by tile
6. benchmark shape sweep
7. objdump or compiler assembly evidence
8. perf stat or substitute counters
9. bandwidth/roofline gap analysis
10. comparison with oneDNN, OpenBLAS, ggml or llama.cpp
\end{lstlisting}

固定两个 shape。小 shape \code{T=2,K=8,N=6} 用来画原始 layout 到 packed layout 的地址变化，并手工验证输出；大 shape \code{H=4096,K=4096,O=4096} 用来估算权重读流、带宽上限、cache/TLB 压力和实际 tokens/s。这样读者既不会被大模型数字淹没，也不会停在没有工程意义的小例子。

\section*{Runtime/serving 的硬交付}

第三册必须明确：AI Infra 不只是 kernel。一个本地推理服务至少要有下面对象：

\begin{lstlisting}[numbers=none,caption={serving runtime 最小对象}]
SessionState:
  request id, token buffer, position, KV cache, sampler state

AdmissionController:
  context limit, KV budget, workspace, queue capacity

Scheduler:
  prefill queue, decode ready queue, cancellation, backpressure

StreamEvent:
  token, final, error, trace, finish reason

ServiceSloReport:
  queue wait, TTFT, TPOT, p95, rejection rate, cancel reclaim
\end{lstlisting}

服务化章节要使用固定 workload：短 prompt、长 RAG prompt、取消请求、超长拒绝、慢客户端 backpressure、后端 fallback。每个 workload 都要输出 trace。否则“能做业务开发”只是说法，不能证明。

\section*{GPU 和分布式的边界}

第三册当前主线是本地推理和 CPU-first 后端，因此 GPU 和分布式不能空喊“也支持”。它们的边界要写清：

\begin{itemize}
  \item GPU 最低要求是 device residency、stream/event、workspace、copy/sync、fallback、StepTrace；不要求一开始写完整 CUDA 生产 kernel，但必须知道端到端延迟不等于 kernel 时间。
  \item 分布式推理最低要求是说明张量并行、pipeline 并行、KV cache 分片、跨节点通信、失败和 SLO 会引入什么新合同；若不实现，就必须在路线图中标为未覆盖，而不是暗示已经具备。
  \item serving 最低要求是单机多 session 的 admission、调度、取消、backpressure 和 SLO；多机 serving 是下一阶段，不应混进当前证据。
\end{itemize}

这种边界不是降标准，而是防止作品集虚胖。强团队更看重你能准确说明“已经证明什么、还没证明什么、下一步怎么证明”。

\section*{工业对标的最低标准}

第三册必须对成熟系统保持敬畏。对标不是只报一个 tokens/s，而是拆阶段：

\begin{center}
\begin{tabular}{p{0.18\linewidth}p{0.36\linewidth}p{0.30\linewidth}}
\toprule
阶段 & 本项目应报告 & 对标时要问 \\
\midrule
load/prepare & mmap/load、packing、device upload & 成熟框架是否预处理或缓存权重 \\
prefill & prompt tokens/s、TTFT 拆分 & 是否使用高质量 GEMM/attention kernel \\
decode & TPOT、p50/p95、context bucket & KV layout、batch、thread/GPU 策略差异 \\
memory & weights、packed、KV、arena peak & 是否使用 paging、quant KV、共享 prefix \\
quality & logits/top-k/sequence drift & tokenizer、quant、sampler 是否一致 \\
\bottomrule
\end{tabular}
\end{center}

至少选择 oneDNN/OpenBLAS/ggml/llama.cpp/ONNX Runtime 中的两个做拆项对照。对照结果落后并不可怕；无法解释差距才可怕。

本册重新分为六个大章。

\topic{第一章：总架构与业务闭环。} 先解释编译器到底在做什么，再把传统编译流水线映射到 AI 推理引擎，并明确第三册怎样继承第一册的 ABI/汇编/测量能力和第二册的硬件/并发/运行时能力。这里会定义全书五层架构：模型资产层、编译前端层、IR 与优化层、后端层、运行时层，并用一条端到端链路把模型文件、manifest、typed graph、executable plan、prefill、KV cache、decode、sampler、oracle 和 profiler 串起来。本册还会明确贯穿项目：一个可解释的 C++20 本地推理服务原型，包含 CLI、服务接口、报告文件、测试夹具、reference、oracle、profiler、资源预算、业务状态机、队列、调度、取消和 backpressure。读者先理解“语义合同”和“机器账本”，后续内容才不会把 manifest、IR、kernel 和 runtime 混成一团。

\topic{第二章：张量、算子、量化与计算账本。} 从真实本地推理引擎边界开始，逐步进入张量布局、地址公式、矩阵向量乘、GEMM、packing、SIMD、量化合同和 7B 计算账本。这里重点回答：一个算子在数学上定义什么，在内存里怎样存在，在机器上怎样读写，低比特表示怎样保持数值合同，模型级 FLOPs、权重字节和 KV cache 容量怎样算清。

\topic{第三章：Transformer 推理、KV cache 与采样。} 把 decoder-only Transformer、reference decoder、prefill、decode、KV cache、logits、sampler 和长上下文状态接起来。这里重点回答：一次真实请求怎样从 prompt 变成 token 流，RoPE、attention、GQA、KV cache、采样策略和停止条件各自改变什么状态，哪些错误必须由 oracle 定位。

\topic{第四章：AI 编译器与推理运行时。} 讲推理引擎里为什么需要编译器，怎样从模型文件导入 typed graph，怎样建立 manifest、schema resolution、shape/dtype/quant verifier，怎样做 memory planning，怎样分层使用 graph IR、tensor IR、loop IR，怎样设计 analysis pass、transform pass、backend lowering 和 executable plan。所有编译原理概念都落到 AI infra 对象上解释。

\topic{第五章：C++20 CPU/GPU 后端优化。} 先建立后端优化总路线：reference、scalar baseline、layout、packing、SIMD、线程、NUMA/TLB、runtime integration、oracle/profiler report。CPU 后端要细到 RAII 所有权、\code{std::span} 视图、tensor view、aligned allocator、packed weight cache、SIMD micro-kernel、寄存器账本、tail path、线程池、NUMA/TLB、KV cache 读流、perf counter、benchmark harness 和成熟框架对标；GPU 后端要从硬件事实落到 device buffer、stream、event、workspace、kernel launch、host-device staging、量化格式支持、fallback 和端到端 latency。

\topic{第六章：工程验收、路线图与作品集证据。} 把 reference、oracle、profiler schema、测试矩阵、覆盖率、性能回归门禁、诊断质量、实施路线图和框架对标收束成作品集证据。这里重点回答：一个学习项目怎样从“能跑”升级为能被强团队认可的 AI Infra 资产。

因此，本册每个大章内部不再堆小标题，而按同一条工程推导线展开：概念定义、最小实现、性能瓶颈、优化设计、实验数据、工业对比、局限复盘。每个优化都要满足同一条验收线：它保持了哪条语义合同，改变了哪条地址流或生命周期，目标后端为什么因此更快，oracle 怎样证明没算错，profiler 怎样证明瓶颈真的下降。只要这些问题没有讲清，就不算完成一个高质量 AI infra 优化。
